% GENERATED FILE. Edit canonical lessons, then run npm run generate:books.
% canonical-source-hash: fnv1a64:ce8576b63af108d0
% canonical-lessons: HI-C545-jua, HI-C545-pasa, HI-C545-bazi, HI-C545-mohra, HI-C545-mat, HI-R545-review-538-545

\chapter{A Game of Chance}
\label{ch:hi-a-game-of-chance}
\hlchaptermodality{545}

\begin{chapteropening}
\textbf{By the end of this chapter:} \emph{I can talk about dice, games and winning a round.}

Complete the last lesson of chapter 545: I can talk about dice, games and winning a round.
\end{chapteropening}

\section[\texorpdfstring{\hi{जुआ} (juā)}{जुआ (juā)}]{\hi{जुआ} (juā) — gambling}

\label{lesson:HI-C545-jua}



\begin{quote}
Before the new one: say the Hindi for an opportunity (formal), then the Hindi for a turn (in a game).
\end{quote}

\subsection*{You'll want to know: \hi{जुआ}}
\textbf{\hi{जुआ}} — \emph{juā} — \textquotedblleft{}gambling\textquotedblright{}.

\begin{quote}
\hi{जुआ} \hi{खेलना} \hi{एक} \hi{बुरी} \hi{आदत} \hi{है।}
\end{quote}

\begin{quote}
\emph{Juā khelnā ek burī ādat hai.}
\end{quote}

\begin{quote}
Gambling is a bad habit.
\end{quote}

Masculine.

\begin{cousinweb}
\textbf{\hi{जुआ}} (\emph{juā}) is a tadbhava, inherited from Sanskrit \emph{dyūta}, \textquotedblleft{}gambling\textquotedblright{} (CDIAL 6608).
\end{cousinweb}

\begin{grammarlens}[title={What you've built}]
One more word for saying what happens if.
\end{grammarlens}

\subsection*{Your turn}
\begin{itemize}
\raggedright
  \item \emph{Say it:} \emph{juā}
  \item \emph{Say it:} \emph{juā}, once more
  \item \emph{Say it:} \emph{Juā khelnā ek burī ādat hai.} — Gambling is a bad habit.
  \item \emph{Recall:} say the Hindi for an opportunity (formal), then the Hindi for a turn (in a game), then say \emph{juā} again
\end{itemize}

\begin{tcolorbox}[breakable,colback=teal!4,colframe=teal!35!black,title={Before you move on}]
What does \hi{जुआ} mean? (\textquotedblleft{}Gambling\textquotedblright{}.) Say it once more.
\end{tcolorbox}
\section[\texorpdfstring{\hi{पासा} (pāsā)}{पासा (pāsā)}]{\hi{पासा} (pāsā) — a die; dice}

\label{lesson:HI-C545-pasa}



\begin{quote}
Before the new one: say the Hindi for a turn (in a game), then the Hindi for gambling.
\end{quote}

\subsection*{You'll want to know: \hi{पासा}}
\textbf{\hi{पासा}} — \emph{pāsā} — \textquotedblleft{}a die; dice\textquotedblright{}.

\begin{quote}
\hi{उसने} \hi{पासा} \hi{फेंका} \hi{और} \hi{छह} \hi{आया।}
\end{quote}

\begin{quote}
\emph{Usne pāsā pheṅkā aur chah āyā.}
\end{quote}

\begin{quote}
He threw the die and got a six.
\end{quote}

Masculine.

\begin{cousinweb}
\textbf{\hi{पासा}} (\emph{pāsā}) is a tadbhava, inherited from Sanskrit \emph{pāśa}, \textquotedblleft{}a die, dice\textquotedblright{} (CDIAL 8132).
\end{cousinweb}

\begin{grammarlens}[title={What you've built}]
One more word for saying what happens if.
\end{grammarlens}

\subsection*{Your turn}
\begin{itemize}
\raggedright
  \item \emph{Say it:} \emph{pāsā}
  \item \emph{Say it:} \emph{pāsā}, once more
  \item \emph{Say it:} \emph{Usne pāsā pheṅkā aur chah āyā.} — He threw the die and got a six.
  \item \emph{Recall:} say the Hindi for a turn (in a game), then the Hindi for gambling, then say \emph{pāsā} again
\end{itemize}

\begin{tcolorbox}[breakable,colback=teal!4,colframe=teal!35!black,title={Before you move on}]
What does \hi{पासा} mean? (\textquotedblleft{}A die; dice\textquotedblright{}.) Say it once more.
\end{tcolorbox}
\section[\texorpdfstring{\hi{बाज़ी} (bāzī)}{बाज़ी (bāzī)}]{\hi{बाज़ी} (bāzī) — a round (of a game); a bet}

\label{lesson:HI-C545-bazi}



\begin{quote}
Before the new one: say the Hindi for gambling, then the Hindi for a die.
\end{quote}

\subsection*{You'll want to know: \hi{बाज़ी}}
\textbf{\hi{बाज़ी}} — \emph{bāzī} — \textquotedblleft{}a round (of a game); a bet\textquotedblright{}.

\begin{quote}
\hi{खेल} \hi{की} \hi{पहली} \hi{बाज़ी} \hi{मीरा} \hi{ने} \hi{जीती।}
\end{quote}

\begin{quote}
\emph{Khel kī pahlī bāzī Mīrā ne jītī.}
\end{quote}

\begin{quote}
Mira won the first round of the game.
\end{quote}

Feminine.

\begin{cousinweb}
\textbf{\hi{बाज़ी}} (\emph{bāzī}) is usually traced to Persian \emph{bāzī}, \textquotedblleft{}play, a game\textquotedblright{}.
\end{cousinweb}

\begin{grammarlens}[title={What you've built}]
One more word for saying what happens if.
\end{grammarlens}

\subsection*{Your turn}
\begin{itemize}
\raggedright
  \item \emph{Say it:} \emph{bāzī}
  \item \emph{Say it:} \emph{bāzī}, once more
  \item \emph{Say it:} \emph{Khel kī pahlī bāzī Mīrā ne jītī.} — Mira won the first round of the game.
  \item \emph{Recall:} say the Hindi for gambling, then the Hindi for a die, then say \emph{bāzī} again
\end{itemize}

\begin{tcolorbox}[breakable,colback=teal!4,colframe=teal!35!black,title={Before you move on}]
What does \hi{बाज़ी} mean? (\textquotedblleft{}A round (of a game); a bet\textquotedblright{}.) Say it once more.
\end{tcolorbox}
\section[\texorpdfstring{\hi{मोहरा} (mohrā)}{मोहरा (mohrā)}]{\hi{मोहरा} (mohrā) — a piece (in chess or a board game)}

\label{lesson:HI-C545-mohra}



\begin{quote}
Before the new one: say the Hindi for a die, then the Hindi for a round (of a game).
\end{quote}

\subsection*{You'll want to know: \hi{मोहरा}}
\textbf{\hi{मोहरा}} — \emph{mohrā} — \textquotedblleft{}a piece (in chess or a board game)\textquotedblright{}.

\begin{quote}
\hi{उसने} \hi{अपना} \hi{मोहरा} \hi{आगे} \hi{बढ़ाया।}
\end{quote}

\begin{quote}
\emph{Usne apnā mohrā āge baṛhāyā.}
\end{quote}

\begin{quote}
He moved his piece forward.
\end{quote}

Masculine.

\begin{cousinweb}
\textbf{\hi{मोहरा}} (\emph{mohrā}) is usually traced to Persian \emph{muhra}, \textquotedblleft{}a bead, a playing piece\textquotedblright{}.
\end{cousinweb}

\begin{grammarlens}[title={What you've built}]
One more word for saying what happens if.
\end{grammarlens}

\subsection*{Your turn}
\begin{itemize}
\raggedright
  \item \emph{Say it:} \emph{mohrā}
  \item \emph{Say it:} \emph{mohrā}, once more
  \item \emph{Say it:} \emph{Usne apnā mohrā āge baṛhāyā.} — He moved his piece forward.
  \item \emph{Recall:} say the Hindi for a die, then the Hindi for a round (of a game), then say \emph{mohrā} again
\end{itemize}

\begin{tcolorbox}[breakable,colback=teal!4,colframe=teal!35!black,title={Before you move on}]
What does \hi{मोहरा} mean? (\textquotedblleft{}A piece (in chess or a board game)\textquotedblright{}.) Say it once more.
\end{tcolorbox}
\section[\texorpdfstring{\hi{मात} (māt)}{मात (māt)}]{\hi{मात} (māt) — defeat; checkmate}

\label{lesson:HI-C545-mat}



\begin{quote}
Before the new one: say the Hindi for a round (of a game), then the Hindi for a piece (in chess or a board game).
\end{quote}

\subsection*{You'll want to know: \hi{मात}}
\textbf{\hi{मात}} — \emph{māt} — \textquotedblleft{}defeat; checkmate\textquotedblright{}.

\begin{quote}
\hi{खेल} \hi{में} \hi{मीरा} \hi{ने} \hi{अरुण} \hi{को} \hi{मात} \hi{दी।}
\end{quote}

\begin{quote}
\emph{Khel meṁ Mīrā ne Aruṇ ko māt dī.}
\end{quote}

\begin{quote}
Mira beat Arun at the game.
\end{quote}

Feminine.

\begin{cousinweb}
\textbf{\hi{मात}} (\emph{māt}) is usually traced to Persian \emph{māt}, \textquotedblleft{}checkmated, at a loss\textquotedblright{}.
\end{cousinweb}

\begin{grammarlens}[title={What you've built}]
One more word for saying what happens if.
\end{grammarlens}

\subsection*{Your turn}
\begin{itemize}
\raggedright
  \item \emph{Say it:} \emph{māt}
  \item \emph{Say it:} \emph{māt}, once more
  \item \emph{Say it:} \emph{Khel meṁ Mīrā ne Aruṇ ko māt dī.} — Mira beat Arun at the game.
  \item \emph{Recall:} say the Hindi for a round (of a game), then the Hindi for a piece (in chess or a board game), then say \emph{māt} again
\end{itemize}

\begin{tcolorbox}[breakable,colback=teal!4,colframe=teal!35!black,title={Before you move on}]
What does \hi{मात} mean? (\textquotedblleft{}Defeat; checkmate\textquotedblright{}.) Say it once more.
\end{tcolorbox}
\section[(dialogue)]{Looking back: from Getting Things Ready to A Game of Chance}

\label{lesson:HI-R545-review-538-545}



\begin{quote}
Say the Hindi for to treasure and ready.
\end{quote}

\subsection*{Your turn}
Say these aloud:

\begin{itemize}
\raggedright
  \item to treasure, ready, management, to be decorated, to pick out, to pick over (grain)
  \item a place of pilgrimage, a prayer, a blessing (for someone), a prayer, a wish (for someone's good), to resolve
  \item experience (gained over time), a clever way, a means, a way (of doing something), a trick, an effort, an endeavour (formal)
  \item an obstacle, a hold-up, a hassle, a tiresome business, to get (something) stuck, a noose, to tangle
  \item talks, a concession, agreement, an offer, a mediator
  \item a word given, an oath, willing, an oath (formal), bound (by a rule or a time)
  \item luck, a coincidence, a chance, an opportunity (formal), a turn (in a game)
  \item gambling, a die, a round (of a game), a piece (in chess or a board game), defeat
\end{itemize}

\begin{tcolorbox}[breakable,colback=teal!4,colframe=teal!35!black,title={Before you move on}]
To treasure? (\textbf{\hi{सँजोना}}, \emph{sãjonā}.) Ready? (\textbf{\hi{तैयार}}, \emph{taiyār}.) Management? (\textbf{\hi{प्रबंध}}, \emph{prabandh}.) To be decorated? (\textbf{\hi{सजना}}, \emph{sajnā}.) To pick out, to pick over (grain)? (\textbf{\hi{बीनना}}, \emph{bīnnā}.) A place of pilgrimage? (\textbf{\hi{तीर्थ}}, \emph{tīrth}.) A prayer, a blessing (for someone)? (\textbf{\hi{दुआ}}, \emph{duā}.) A prayer? (\textbf{\hi{प्रार्थना}}, \emph{prārthnā}.) A wish (for someone's good)? (\textbf{\hi{कामना}}, \emph{kāmnā}.) To resolve? (\textbf{\hi{ठानना}}, \emph{ṭhānnā}.) Experience (gained over time)? (\textbf{\hi{तजुर्बा}}, \emph{tajurbā}.) A clever way? (\textbf{\hi{युक्ति}}, \emph{yukti}.) A means, a way (of doing something)? (\textbf{\hi{ज़रिया}}, \emph{zariyā}.) A trick? (\textbf{\hi{तरकीब}}, \emph{tarkīb}.) An effort, an endeavour (formal)? (\textbf{\hi{प्रयत्न}}, \emph{prayatn}.) An obstacle, a hold-up? (\textbf{\hi{रुकावट}}, \emph{rukāvaṭ}.) A hassle, a tiresome business? (\textbf{\hi{झमेला}}, \emph{jhamelā}.) To get (something) stuck? (\textbf{\hi{अटकाना}}, \emph{aṭkānā}.) A noose? (\textbf{\hi{फंदा}}, \emph{phandā}.) To tangle? (\textbf{\hi{उलझाना}}, \emph{uljhānā}.) Talks? (\textbf{\hi{वार्ता}}, \emph{vārtā}.) A concession? (\textbf{\hi{रियायत}}, \emph{riyāyat}.) Agreement? (\textbf{\hi{सहमति}}, \emph{sahmati}.) An offer? (\textbf{\hi{पेशकश}}, \emph{peśkaś}.) A mediator? (\textbf{\hi{मध्यस्थ}}, \emph{madhyasth}.) A word given? (\textbf{\hi{वचन}}, \emph{vacan}.) An oath? (\textbf{\hi{क़सम}}, \emph{qasam}.) Willing? (\textbf{\hi{राज़ी}}, \emph{rāzī}.) An oath (formal)? (\textbf{\hi{शपथ}}, \emph{śapath}.) Bound (by a rule or a time)? (\textbf{\hi{पाबंद}}, \emph{pāband}.) Luck? (\textbf{\hi{क़िस्मत}}, \emph{qismat}.) A coincidence? (\textbf{\hi{संयोग}}, \emph{saṁyog}.) A chance? (\textbf{\hi{मौक़ा}}, \emph{mauqā}.) An opportunity (formal)? (\textbf{\hi{अवसर}}, \emph{avsar}.) A turn (in a game)? (\textbf{\hi{दाँव}}, \emph{dā̃v}.) Gambling? (\textbf{\hi{जुआ}}, \emph{juā}.) A die? (\textbf{\hi{पासा}}, \emph{pāsā}.) A round (of a game)? (\textbf{\hi{बाज़ी}}, \emph{bāzī}.) A piece (in chess or a board game)? (\textbf{\hi{मोहरा}}, \emph{mohrā}.) Defeat? (\textbf{\hi{मात}}, \emph{māt}.)
\end{tcolorbox}
